\begin{problem}{L}{LARPing the Length}{1 second}

\noindent
On a normal Saturday, the police headquarters in Beika received a rather strange call. A self-proclaimed bomber named ``Nanafushi'' claimed to have hidden a bomb somewhere in the city, but refused to reveal its location. Instead, Nanafushi promised to disclose it only if the police could win a peculiar game.

\noindent
Nanafushi claimed that they were, in fact, a giant stick of some secret length $N$. For each query, the police may present two sticks of lengths $A$ and $B$, and Nanafushi will imagine the three sticks --- $A$, $B$, and their own length $N$ --- forming a triangle. If the three sticks cannot form a non-degenerate triangle, Nanafushi becomes furious and immediately ends the call. Otherwise, Nanafushi simply tells the police whether the resulting triangle is acute or not.

\noindent
Three sticks of lengths $x$, $y$, and $z$ can form a non-degenerate triangle iff $x+y+z>2 \max(x,y,z)$ and the triangle is acute iff, in addition, $x^2+y^2+z^2>2 \max(x,y,z)^2$.

\noindent
Your task is to help the police determine the value of $N$ using \textbf{at most 67 queries}.

\section*{\sffamily Interaction}

\noindent
Your program first reads a single integer $t$ ($1 \le t \le 100$) --- the number of test cases.

\noindent
For each test case, the secret stick length $N$ ($2 \le N \le 10^9$) is fixed before the interaction begins.

\noindent
To make a query, print a line in the following format:
\begin{itemize}
\item \texttt{?} $A$ $B$
\end{itemize}

\noindent
where $1 \le A,B \le 10^9$.

\noindent
Nanafushi considers three sticks of lengths $A$, $B$, and $N$. If they cannot form a \textbf{non-degenerate triangle}, the interaction immediately ends and you receive a \texttt{Wrong Answer} verdict.

\noindent
Otherwise, after printing a query, flush the output and read a single integer $s$:
\begin{itemize}
\item $s = 0$ --- if the triangle formed by the three sticks is not acute.
\item $s = 1$ --- if the triangle formed by the three sticks is acute.
\end{itemize}

\noindent
Once you have determined the secret length $N$, print:
\begin{itemize}
\item \texttt{!} $N$
\end{itemize}

\noindent
Printing the answer does \textbf{not} count as a query. You may perform at most \textbf{67 queries} per test case. If your program performs more than $67$ queries, outputs an invalid query, or makes a query whose three sticks cannot form a non-degenerate triangle, you may receive a \texttt{Wrong Answer} verdict.

\noindent
The interactor is \textbf{not adaptive}; the value of $N$ is fixed before the interaction begins and does not depend on your queries.

\noindent
After every query and the final answer, output the end of line and flush the output. Otherwise, you may get a \texttt{Time Limit Exceeded} verdict. To do this, use:

\begin{itemize}
\item \texttt{fflush(stdout)} or \texttt{cout.flush()} in C/C++;
\item \texttt{System.out.flush()} in Java;
\item \texttt{sys.stdout.flush()} in Python.
\end{itemize}

\begin{tabular*}{\textwidth}{@{\extracolsep{\fill}}l c r}
  \sffamily{\textbf{Read}} & \sffamily{\textbf{Sample Interaction}} & \sffamily{\textbf{Write}}
\end{tabular*}

\hfill\iobox{2}\\
\vspace{1mm}
\iobox{? 2 2}\hfill\\
\vspace{1mm}
\hfill\iobox{1}\\
\vspace{1mm}
\iobox{? 2 3}\hfill\\
\vspace{1mm}
\hfill\iobox{0}\\
\vspace{1mm}
\iobox{! 2}\hfill\\
\vspace{1mm}
\iobox{? 3 4}\hfill\\
\vspace{1mm}
\hfill\iobox{0}\\
\vspace{1mm}
\iobox{? 4 4}\hfill\\
\vspace{1mm}
\hfill\iobox{1}\\
\vspace{1mm}
\iobox{! 5}\hfill

\pagebreak

\end{problem}